\section{Profiler}
PDblend 使用离线 profiler 将设备行为转化为在线调度与资源配置可查询的代价模型。
它按模型、张量并行度、频率及请求形状区分数据，分别采集 prefill、decode、静态状态和实例间 KV 交接，输出时延、功率、KV 容量与状态切换耗时。
整个流程分为采集、校准和查询三层：原始测量保留环境与实现身份；校准仅使用训练样本，并在独立验证前冻结参数；在线模块读取已经验证且覆盖当前输入的预测。

\paragraph{采集实际执行形状。}
Prefill 的主要变量是输入长度，decode 则同时依赖 batch size 和上下文长度。
Native 采集从引擎内部读取 CUDA 执行事件，并核对各 TP rank 的请求和调度形状，以最慢 rank 的耗时作为一次逻辑执行的耗时。
Decode 只在全部请求完成 prefill、进入连续生成后测量；每个功率点包含至少三次独立重复，每次先稳定至少两秒，再采样至少五秒。
上下文随生成增长，因此模型使用窗口内实际执行的上下文，而非直接使用名义 prompt 长度。
Mixed 探针另外观测插入 prefill 时造成的 decode stall，用于检查两种阶段共置时的干扰。

\paragraph{有界时延与功率模型。}
令 $n$ 为 prefill token 数，$B$ 为 decode batch，$c$ 为上下文长度，$f$ 为 GPU 频率。
当前 native timing 采用非负系数拟合：
\begin{align}
T_p(n,f)&=a_f+b_fn+d_fn^2,\label{eq:profile-prefill}\\
T_d(B,c,f)&=u_f+v_fB+w_fBc.\label{eq:profile-decode}
\end{align}
代码中的尺度归一化与单位换算已吸收进系数。
Prefill 允许在实测长度区间内查询，decode 进一步要求 $(B,c)$ 位于训练形状的覆盖域；频率必须属于已经测量的档位。
功率由同一实例各 GPU 的采样积分得到：
\begin{equation}
\bar P=\frac{1}{t_1-t_0}\int_{t_0}^{t_1}\sum_{g\in G}P_g(t)\,dt,
\qquad E=\bar P(t_1-t_0).
\label{eq:profile-energy}
\end{equation}
Decode 功率保存在有界表中，随实际上下文插值；低 batch $1,2,3$ 仅接受分别验证的精确节点。
请求周期的平均功率与 CUDA 内核耗时具有不同边界，不能直接相乘；完整服务能量需要对应周期或布局的独立模型。

\paragraph{验证与运行时接口。}
独立 holdout 检查预测误差、重复稳定性、实际频率和样本完整性，且不能回流参与拟合。
发布时绑定原始样本、冻结参数和验证记录，使同一份代价模型能够离线重放，避免不同实现或测量口径的结果被混用。
在线查询同时检查模型身份和覆盖范围，缺失的频率、形状或运行时组件会使相应候选不可用。
KV 交接以实际 P/D 请求的输出时序标定，状态组件则提供静态功率、唤醒和调频耗时；这些量随后进入请求路径和池配置的可行性判断。
历史 profile 仍保留兼容接口，但单个时延或功率组件通过验证，并不自动意味着整个在线系统已获得完整能量与切换行为的验证。
